% !TEX encoding = UTF-8 Unicode
% !TEX spellcheck = en_US

\documentclass[a4paper, 11pt]{article}

\RequirePackage[l2tabu, orthodox]{nag}
\usepackage[utf8]{inputenc}
\usepackage[T1]{fontenc}
\usepackage{lmodern}
\usepackage[margin=2cm]{geometry}
\usepackage{enumitem}
\usepackage{fancyhdr}
\setlength{\headheight}{14pt}
\usepackage{lastpage}
\usepackage{hyperref}

% Specific commands
\newcommand{\FullSWOFnD}{\emph{FullSWOF\_1D/2D}}
\newcommand{\FullSWOFoneD}{\emph{FullSWOF\_1D}}
\newcommand{\FullSWOFtwoD}{\emph{FullSWOF\_2D}}
\newcommand{\SWASHES}{\emph{SWASHES}}
\newcommand{\contactFullName}{Frédéric \textsc{Darboux}}
\newcommand{\contactEmail}{\href{mailto:Frederic.Darboux@inrae.fr}{Frederic.Darboux@inrae.fr}}
\newcommand{\MainWebSiteFSoneD}{https://sourcesup.renater.fr/projects/fullswof-1d/}
\newcommand{\MainWebSiteFStwoD}{https://sourcesup.renater.fr/projects/fullswof-2d/}
\newcommand{\MainWebSiteSWASHES}{https://sourcesup.renater.fr/projects/swashes/}

% Headers and footers
\lhead{Running \FullSWOFnD{} or \SWASHES{} under windows}
\rhead{INRAE-IGE --- Univ. Orléans-IDP}
\cfoot{\thepage /\pageref{LastPage}}
\pagestyle{fancy}

% Lists: no spacing between items
\setlist{noitemsep}

\title{Application note:\\
Using Cygwin to compile and run\\
\FullSWOFoneD, \FullSWOFtwoD{} or \SWASHES{}\\
under windows.}

\author{\contactFullName, \contactEmail}
\date{2023-03-20}

\begin{document}

\maketitle

\thispagestyle{fancy}

\FullSWOFoneD, \FullSWOFtwoD{} and \SWASHES{} have been developed under Unix-like environments. Although it may not be required, it is convenient to use such a Unix-like environment under windows. This application note gives directions about installing the Unix-like environment Cygwin, and using it to compile and run  
\href{\MainWebSiteFSoneD}{\FullSWOFoneD},
\href{\MainWebSiteFStwoD}{\FullSWOFtwoD}, and
\href{\MainWebSiteSWASHES}{\SWASHES}.

\section{Installation of Cygwin}

\begin{enumerate}
\item From \url{www.cygwin.com}, download the file \verb!setup-x86_64.exe!, and save it into a dedicated directory (e.g. \verb!c:\software\cygwin_install!).
\item Launch \verb!setup-x86_64.exe!.
\item Click on the ``Next'' button on the first screens. Then, choose a download site (next to your location), and click on the ``Next'' button once more. You will be prompted to select packages. A basic set of packages is already selected. To compile the software smoothly, you need to add a few more packages. For the following list of packages, in the column entitled ``New``, click on the drop-down menu to select the version to be installed:
	\begin{itemize}
	\item in the category ``Archive''
		\begin{itemize}
		\item unzip: Info-ZIP decompression utility
		\end{itemize}
	\item in the category ``Devel''
		\begin{itemize}
		\item gcc-g++: C++ compiler Collection (C++)
		\item make: The GNU version of the 'make' utility
		\end{itemize}
	\end{itemize}
Additionally, if you want to display graphs using gnuplot, you also need:
	\begin{itemize}
	\item in the category ``Graphics''
		\begin{itemize}
		\item gnuplot-X11: A command-line driven interactive function plotting utility
		\end{itemize}
	\item in the category ``X11''
		\begin{itemize}
		\item xorg-server: X.org X server
		\item xinit: X.org X server launcher
		\end{itemize}
	\end{itemize}
\item Then click ``Next'' up to the start of downloading.
\item Finally, click ``End''.
\end{enumerate}

This will use about of 610~MB (840~MB with gnuplot).

\section{Compiling and running the software}

\begin{enumerate}
\item Download the \FullSWOFoneD\footnote{\url{\MainWebSiteFSoneD}}, \FullSWOFtwoD{}\footnote{\url{\MainWebSiteFStwoD}} or \SWASHES{}\footnote{\url{\MainWebSiteSWASHES}} package into a dedicated directory.
\item Open a Cygwin terminal using the desktop icon or the menu.
\item Move to your dedicated directory. For example, to access the directory \verb!D:\user\code!, you should enter the command \verb!cd /cygdrive/d/user/code!
\item Unzip the package (e.g. \verb!unzip package.zip!).
\item Move to the newly-created directory. This directory will contain all the files related to the source code.
\item Refer to the software-specific documentation to compile and run the code.
\item To use gnuplot, launch the ``XWin Server'' from the windows' menu. Then, in the taskbar, select \verb!Xserver -> System tools -> Cygwin terminal!. In the cygwin terminal, type in \verb!gnuplot!. To test it, type \verb!test! after the gnuplot prompt. An example window should display.

\end{enumerate}

\emph{For more information about Cygwin, see \url{http://www.cygwin.com}}

\end{document}
